% Chapitre 2 - Page 7 | Travaux Avant-Projet | L3 ELN
\documentclass[11pt,a4paper]{article}
\usepackage{fontspec}
\setmainfont{TeXGyrePagella}
\usepackage[french]{babel}
\usepackage{graphicx,tikz,xcolor,geometry,hyperref}
\usetikzlibrary{calc}
\geometry{margin=0pt}
\pagestyle{empty}
\definecolor{bleufonce}{HTML}{1B3A6B}
\definecolor{bleugris}{HTML}{5C7190}
\hypersetup{colorlinks=true,urlcolor=bleufonce}
\newcommand{\pos}[4]{\node[anchor=north west,inner sep=0pt,outer sep=0pt,text width=#3,align=left] at ($(current page.north west)+(#1,-#2)$) {#4};}
\newcommand{\sect}[1]{{\color{bleufonce}\bfseries\fontsize{17}{19}\selectfont #1}}
\newcommand{\step}[1]{{\color{bleufonce}\bfseries\fontsize{12.5}{14.5}\selectfont #1}}
\newcommand{\smallfont}{\fontsize{9.8}{12}\selectfont}
\newcommand{\legende}[1]{{\color{bleugris}\fontsize{8.5}{10}\selectfont #1}}
\newcommand{\puce}{\textcolor{bleufonce}{\textbullet}\enspace}
% \shot{x}{y}{largeur}{hauteur}{fichier} : capture encadrée
\newcommand{\shot}[5]{\node[anchor=north west,inner sep=0pt,outer sep=0pt] at ($(current page.north west)+(#1,-#2)$) {\includegraphics[width=#3]{#5}};
  \draw[bleufonce!75,line width=.55pt,rounded corners=1mm]
    ($(current page.north west)+(#1,-#2)+(-.4mm,.4mm)$) rectangle ($(current page.north west)+(#1,-#2)+(#3,-#4)+(.4mm,-.4mm)$);}
\begin{document}
\begin{tikzpicture}[remember picture,overlay]
\draw[bleufonce,line width=.95pt,rounded corners=8mm]
 ($(current page.north west)+(5mm,-5mm)$) rectangle ($(current page.south east)+(-5mm,5mm)$);

\pos{10mm}{13mm}{189mm}{\sect{5. Créer un circuit dans Tinkercad}}
\pos{10mm}{25mm}{188mm}{{\fontsize{10.5}{13}\selectfont Les trois étapes ci-dessous sont illustrées par des captures du site officiel \href{https://www.tinkercad.com/circuits}{\textbf{www.tinkercad.com/circuits}}.}}

% --- Étape 1 ---------------------------------------------------------------------
\pos{10mm}{37mm}{188mm}{\step{Étape 1 --- Créer un compte ou se connecter}}
\shot{14mm}{46mm}{40mm}{55.7mm}{assets/T02_compte.png}
\pos{62mm}{47mm}{135mm}{{\smallfont Sur le site officiel, cliquer sur \textbf{S'inscrire} lors de la première visite, ou sur \textbf{Connexion} si le compte existe déjà. Tinkercad demande ensuite le type de compte :\\[2mm]
  \puce \textbf{Étudiants disposant d'un code de classe} : l'enseignant a créé une classe et donne un code ;\\[1.3mm]
  \puce \textbf{Comptes d'étudiants} : compte individuel pour l'apprentissage ;\\[1.3mm]
  \puce \textbf{Enseignants} : pour créer et suivre des classes ;\\[1.3mm]
  \puce \textbf{Comptes personnels} : usage personnel, en dehors d'un cadre scolaire.}}
\pos{14mm}{104mm}{184mm}{\legende{Figure 6 --- Choix du type de compte lors de l'inscription.}}

% --- Étape 2 ---------------------------------------------------------------------
\pos{10mm}{114mm}{188mm}{\step{Étape 2 --- Ouvrir l'espace Circuits}}
\shot{14mm}{123mm}{124mm}{55.9mm}{assets/T03_tableau_bord.png}
\pos{145mm}{124mm}{52mm}{{\smallfont Une fois connecté, ouvrir le menu \textbf{Concevoir}. Il présente les espaces de Tinkercad : \textbf{Conception 3D}, \textbf{Circuits}, \textbf{Codeblocks}\ldots\\[2mm]
  Choisir \textbf{Circuits} pour accéder aux montages électroniques.}}
\pos{14mm}{181mm}{184mm}{\legende{Figure 7 --- Le menu \textbf{Concevoir} et l'espace \textbf{Circuits}.}}

% --- Étape 3 ---------------------------------------------------------------------
\pos{10mm}{191mm}{188mm}{\step{Étape 3 --- Créer un nouveau circuit}}
\shot{14mm}{200mm}{44mm}{43.7mm}{assets/T04_creer.png}
\pos{66mm}{201mm}{131mm}{{\smallfont Cliquer sur le bouton \textbf{+ Créer}, puis choisir \textbf{Circuits}.\\[2mm]
  Un plan de travail vide s'ouvre : c'est l'interface présentée à la page précédente (Figure 5). Le nom du projet, en haut à gauche, peut être modifié d'un simple clic.}}
\pos{14mm}{246mm}{184mm}{\legende{Figure 8 --- Le bouton \textbf{+ Créer} et le choix \textbf{Circuits}.}}

\draw[bleufonce,line width=.7pt,rounded corners=2mm]
 ($(current page.north west)+(10mm,-255mm)$) rectangle ($(current page.north west)+(200mm,-270mm)$);
\pos{15mm}{257.3mm}{180mm}{{\fontsize{9.4}{11.4}\selectfont \textcolor{bleufonce}{\bfseries À retenir :} un compte suffit pour accéder à Tinkercad depuis n'importe quel ordinateur ; les circuits sont enregistrés automatiquement.}}

\draw[bleufonce!55,line width=.45pt] ($(current page.south west)+(10mm,14.7mm)$)--($(current page.south east)+(-10mm,14.7mm)$);
\node[anchor=south west,inner sep=0pt,text=bleufonce!80] at ($(current page.south west)+(10mm,7.4mm)$) {\fontsize{9.5}{11}\selectfont 2026/2027};
\node[anchor=south east,inner sep=0pt,text=bleufonce!80] at ($(current page.south east)+(-10mm,7.4mm)$) {\fontsize{9.5}{11}\selectfont Page 7};
\end{tikzpicture}
\null
\end{document}
